%!TEX TS-program = xelatex
\documentclass[a4paper,12pt]{article}
\usepackage[margin=2.5cm]{geometry}
\usepackage{fontspec}
\usepackage{polyglossia}
\setmainlanguage{french}
\usepackage{xcolor}
\usepackage{amsmath,amssymb}
\usepackage{tikz}
\usetikzlibrary{shapes.geometric, arrows.meta, positioning, calc}
\usepackage{pgfplots}
\pgfplotsset{compat=1.18}
\usepackage{listings}
\usepackage{caption}
\usepackage{subcaption}
\usepackage{hyperref}
\hypersetup{
    colorlinks=true,
    linkcolor=blue,
    urlcolor=blue,
    citecolor=blue
}

% Style TikZ
\tikzstyle{startstop} = [rectangle, rounded corners, minimum width=3cm, minimum height=1cm, text centered, draw=black, fill=red!30]
\tikzstyle{process} = [rectangle, minimum width=3cm, minimum height=1cm, text centered, draw=black, fill=blue!20]
\tikzstyle{decision} = [diamond, minimum width=3cm, minimum height=1cm, text centered, draw=black, fill=green!20]
\tikzstyle{arrow} = [thick,->,>=stealth]

% Style listings
\lstdefinestyle{python}{
    language=Python,
    basicstyle=\ttfamily\small,
    keywordstyle=\color{blue},
    commentstyle=\color{gray},
    stringstyle=\color{red},
    showstringspaces=false,
    numbers=left,
    numberstyle=\tiny\color{gray},
    breaklines=true,
    frame=single,
    captionpos=b
}

\title{\textbf{Analyse du flux de données} \\ de l'optimisation métaleuristique aux graphiques}
\author{}
\date{}

\begin{document}

\maketitle

\tableofcontents
\newpage

\section{Introduction}
Ce document décrit le cheminement des données depuis la recherche métaleuristique de minima du potentiel scalaire $V(s,\tau,\text{coeffs})$ jusqu'à la production des graphiques comparatifs. L'objectif est de comprendre comment les valeurs numériques obtenues par l'algorithme génétique sont ensuite raffinées, classifiées et finalement visualisées.

Le potentiel considéré est issu de compactifications de théorie des cordes et s'écrit, dans sa version complète avec terme de \textit{lifting} $AD5$ :
\[
V_{\text{complète}} = \frac{AH3\,s}{\tau^3} + \frac{AF3}{s\,\tau^3} + \frac{AF5}{\tau^4} + \frac{AO3}{\tau^3} + \frac{AD5}{\sqrt{s}\,\tau^{5/2}}.
\]
La version sans $AD5$ (notée $V_{\text{réduite}}$) est obtenue en annulant ce coefficient.

L'optimisation est réalisée par un algorithme génétique (implémenté avec la bibliothèque \texttt{mealpy}) qui minimise une fonction de \textit{fitness} pénalisant :
\begin{itemize}
    \item un gradient non nul,
    \item une valeur propre de la matrice hessienne négative (tachyon),
    \item une trace négative,
    \item un potentiel négatif (pour les recherches de dS).
\end{itemize}
Les solutions candidates sont ensuite raffinées numériquement par \texttt{FindRoot} (dans Mathematica) ou par une méthode équivalente en Python, puis classifiées en dS stable ($V>0$, valeurs propres $>0$) ou AdS stable ($V<0$, valeurs propres $>0$).

Enfin, pour une solution particulière, on compare le potentiel complet et le potentiel réduit en traçant $V$ en fonction de $s$ (à $\tau$ fixé) et en fonction de $\tau$ (à $s$ fixé).

\section{Organisation des scripts Python}
Deux fichiers principaux sont utilisés :

\begin{description}
    \item[\texttt{MetastableVacua.py}] : contient la définition symbolique du potentiel $V$, de la matrice hessienne, des valeurs propres, ainsi que les fonctions d'évaluation numérique (comme \texttt{potencial\_eval}, \texttt{autovalores\_VHess}) et les fonctions de fitness.
    \item[\texttt{BusquedaGeneticos.py}] : lance l'optimisation métaleuristique (algorithme génétique) pour différentes configurations (dS, AdS, avec ou sans variables fixes). Les résultats sont sauvegardés dans des fichiers CSV.
    \item[\texttt{test\_fitness\_function.py}] : contient des fonctions de validation et de post‑traitement, notamment la lecture des CSV exportés depuis Mathematica et des graphiques de vérification.
\end{description}

\section{Flux des données : de la recherche aux graphiques}
La figure~\ref{fig:flowchart} résume les étapes principales.

\begin{figure}[h]
\centering
\begin{tikzpicture}[node distance=2cm]
\node (start) [startstop] {Définition du potentiel $V$};
\node (ga) [process, below of=start] {Algorithme génétique (\texttt{mealpy})};
\node (csv) [process, below of=ga] {Sauvegarde des solutions brutes (CSV)};
\node (refine) [process, below of=csv] {Raffinement par \texttt{FindRoot} (Mathematica)};
\node (datag) [process, below of=refine] {Solutions raffinées (\texttt{datag})};
\node (class) [decision, below of=datag, yshift=-1cm] {Classification stable ?};
\node (dS) [process, below left of=class, xshift=-2cm] {dS stable};
\node (AdS) [process, below right of=class, xshift=2cm] {AdS stable};
\node (inter) [process, below of=class, yshift=-2cm] {Choix d'un indice \texttt{inter}};
\node (reduced) [process, below of=inter] {Résolution exacte du système réduit ($AD5=0$)};
\node (plot) [startstop, below of=reduced] {Production des graphiques comparatifs};

\draw [arrow] (start) -- (ga);
\draw [arrow] (ga) -- (csv);
\draw [arrow] (csv) -- (refine);
\draw [arrow] (refine) -- (datag);
\draw [arrow] (datag) -- (class);
\draw [arrow] (class) -- node[anchor=east] {oui} (dS);
\draw [arrow] (class) -- node[anchor=west] {oui} (AdS);
\draw [arrow] (dS) |- (inter);
\draw [arrow] (AdS) |- (inter);
\draw [arrow] (inter) -- (reduced);
\draw [arrow] (reduced) -- (plot);
\end{tikzpicture}
\caption{Diagramme de flux des données.}
\label{fig:flowchart}
\end{figure}

\subsection{Recherche métaleuristique}
La fonction \texttt{dS\_search} (ou \texttt{AdS\_search}) dans \texttt{BusquedaGeneticos.py} exécute 100 itérations indépendantes de l'algorithme génétique. Chaque itération produit une solution candidate sous forme d'un vecteur contenant les coefficients $(AH3, AF3, AF5, AO3, AD5)$ et les champs $(s,\tau)$. La valeur de la fonction de fitness (somme des pénalités) est également enregistrée.

Les données collectées pour chaque itération sont :
\begin{itemize}
    \item $V$ (potentiel évalué),
    \item \texttt{fitness},
    \item les coefficients,
    \item $s$, $\tau$,
    \item les deux valeurs propres de la hessienne $\lambda_1,\lambda_2$,
    \item un indicateur « tachyon » si $\min(\lambda_i) < 0$.
\end{itemize}
Ces informations sont sauvegardées dans un fichier CSV (par exemple \texttt{found\_solutions\_\_dS.csv}).

\subsection{Raffinement des solutions}
Les solutions brutes sont ensuite raffinées en résolvant $\nabla V = 0$ avec une grande précision. Dans le notebook Mathematica, cela est fait par \texttt{FindRoot}. Les solutions raffinées sont stockées dans une structure \texttt{datag} (une liste de règles au format Mathematica). Ce raffinement élimine les solutions où le gradient est trop grand et améliore la précision des valeurs propres.

\subsection{Classification}
Chaque solution raffinée est classée selon le signe de $V$ et des valeurs propres :
\begin{itemize}
    \item \textbf{dS stable} : $V > 0$, $\lambda_1 > 0$, $\lambda_2 > 0$,
    \item \textbf{AdS stable} : $V < 0$, $\lambda_1 > 0$, $\lambda_2 > 0$.
\end{itemize}
Les solutions instables (au moins une valeur propre négative) sont écartées des analyses ultérieures.

\subsection{Choix d'une solution et résolution exacte du système réduit}
On sélectionne un indice \texttt{inter} correspondant à une solution d'intérêt (par exemple la sixième). Pour cette solution, on définit un potentiel réduit $V_{\text{réduit}}$ en fixant $AD5 = 0$. On résout alors le système $\partial_s V_{\text{réduit}} = 0$, $\partial_\tau V_{\text{réduit}} = 0$ de manière exacte (via \texttt{solve} de SymPy). On obtient deux solutions analytiques ; on conserve la seconde (\texttt{sol[[2]]}).

\subsection{Production des graphiques}
Pour la solution choisie, on génère deux graphiques :
\begin{enumerate}
    \item \textbf{Graphique en $s$} : on fixe $\tau$ à la valeur de la solution raffinée. On trace deux courbes :
    \begin{itemize}
        \item \textcolor{blue}{bleue} : $V_{\text{complète}}(s)$ (avec tous les coefficients),
        \item \textcolor{red}{rouge} : $V_{\text{réduit}}(s)$ (avec $AD5=0$) en utilisant la solution exacte \texttt{sol[[2]]} pour $(s,\tau)$.
    \end{itemize}
    La plage en $s$ est $[0.5\,s_0,\;2.5\,s_0]$.
    \item \textbf{Graphique en $\tau$} : on fixe $s$ à la valeur de la solution raffinée. On trace :
    \begin{itemize}
        \item \textcolor{blue}{bleue} : $V_{\text{complète}}(\tau)$,
        \item \textcolor{red}{rouge} : $V_{\text{réduit}}(\tau)$ en utilisant cette fois la première solution exacte \texttt{sol[[1]]}.
    \end{itemize}
    La plage en $\tau$ est $[0.5\,\tau_0,\;5.5\,\tau_0]$.
\end{enumerate}
Ces graphiques permettent de visualiser l'accord entre le minimum trouvé par la recherche métaleuristique et la solution exacte du système simplifié, validant ainsi la cohérence de la procédure.

\section{Exemples de graphiques produits}
Les figures ci‑dessous illustrent le type de résultats obtenus.

\begin{figure}[h]
\centering
\begin{subfigure}{0.45\textwidth}
\centering
\begin{tikzpicture}
\begin{axis}[
    width=\textwidth,
    xlabel={$s$},
    ylabel={$V(s,\tau_0)$},
    title={$V(s)$ avec $\tau$ fixé},
    domain=1.5:3.5,
    samples=100,
    grid=both,
    legend pos=north west,
]
\addplot[blue, thick] {0.0002*(x-2.5)^2 - 0.001};
\addplot[red, thick] {0.00025*(x-2.5)^2 - 0.0012};
\legend{$V_{\text{complète}}$, $V_{\text{réduit}}$}
\end{axis}
\end{tikzpicture}
\caption{Comparaison en $s$ (exemple stylisé).}
\label{fig:plot_s}
\end{subfigure}
\hfill
\begin{subfigure}{0.45\textwidth}
\centering
\begin{tikzpicture}
\begin{axis}[
    width=\textwidth,
    xlabel={$\tau$},
    ylabel={$V(s_0,\tau)$},
    title={$V(\tau)$ avec $s$ fixé},
    domain=2:10,
    samples=100,
    grid=both,
    legend pos=north east,
]
\addplot[blue, thick] {0.001*(x-6)^2 - 0.002};
\addplot[red, thick] {0.0012*(x-6)^2 - 0.0025};
\legend{$V_{\text{complète}}$, $V_{\text{réduit}}$}
\end{axis}
\end{tikzpicture}
\caption{Comparaison en $\tau$ (exemple stylisé).}
\label{fig:plot_tau}
\end{subfigure}
\caption{Exemples de graphiques produits après l'analyse.}
\label{fig:plots}
\end{figure}

\section{Code Python pour la génération des graphiques}
La fonction suivante, à intégrer dans un module \texttt{plotting.py}, permet de reproduire les graphiques à partir d'une ligne d'un DataFrame contenant les solutions.

\begin{lstlisting}[style=python, caption={Fonction de tracé pour une solution}]
import numpy as np
import matplotlib.pyplot as plt
from MetastableVacua import potencial_eval, potencial_eval_AdS

def plot_solution_comparison(row, show=True, save=None):
    """
    row : pandas Series contenant AH3, AF3, AF5, AO3, AD5 (optionnel), s, tau
    """
    AH3 = row["AH3"]
    AF3 = row["AF3"]
    AF5 = row["AF5"]
    AO3 = row["AO3"]
    s0 = row["s"]
    tau0 = row["τ"]
    AD5 = row.get("AD5", 0.0)

    # Plages
    s_vals = np.linspace(0.5 * s0, 2.5 * s0, 1000)
    tau_vals = np.linspace(0.5 * tau0, 5.5 * tau0, 1000)

    # Potentiel complet (avec AD5)
    V_s_full = [potencial_eval_AdS(AH3, AF3, AF5, AO3, AD5, tau0, s) for s in s_vals]
    V_tau_full = [potencial_eval_AdS(AH3, AF3, AF5, AO3, AD5, tau, s0) for tau in tau_vals]

    # Potentiel réduit (sans AD5)
    V_s_noad5 = [potencial_eval(AH3, AF3, AF5, AO3, tau0, s) for s in s_vals]
    V_tau_noad5 = [potencial_eval(AH3, AF3, AF5, AO3, tau, s0) for tau in tau_vals]

    fig, (ax1, ax2) = plt.subplots(1, 2, figsize=(12, 5))

    ax1.plot(s_vals, V_s_full, 'b-', label='avec AD5')
    ax1.plot(s_vals, V_s_noad5, 'r-', label='sans AD5')
    ax1.axhline(0, color='k', linestyle=':', alpha=0.5)
    ax1.set_xlabel('s')
    ax1.set_ylabel('V(s, τ₀)')
    ax1.set_title(f'V(s) avec τ = {tau0:.4f}')
    ax1.legend()
    ax1.grid(alpha=0.3)

    ax2.plot(tau_vals, V_tau_full, 'b-', label='avec AD5')
    ax2.plot(tau_vals, V_tau_noad5, 'r-', label='sans AD5')
    ax2.axhline(0, color='k', linestyle=':', alpha=0.5)
    ax2.set_xlabel('τ')
    ax2.set_ylabel('V(s₀, τ)')
    ax2.set_title(f'V(τ) avec s = {s0:.4f}')
    ax2.legend()
    ax2.grid(alpha=0.3)

    plt.tight_layout()
    if save:
        plt.savefig(save, dpi=300)
    if show:
        plt.show()
    return fig
\end{lstlisting}

\section{Conclusion}
Le flux de données présenté permet de passer d'une optimisation stochastique à une validation analytique et visuelle des minima du potentiel. Chaque étape (recherche, raffinement, classification, résolution exacte) est essentielle pour garantir la fiabilité des solutions et la compréhension du paysage du potentiel.

Les graphiques finaux confrontent le potentiel complet (incluant tous les termes) à une version simplifiée où un coefficient est annulé. Ils confirment que la solution trouvée est bien un minimum du système complet et que l'approximation sans $AD5$ en capture correctement les caractéristiques principales autour du point critique.

\end{document}
